\section{Zadanie 3}
Niech \( X \) będzie zmienną Poissona z wartością oczekiwaną \( \mu \in \{1,2,\ldots\} \).

\begin{enumerate}
\item[(i)] Pokaż, że 
\[
\Pr(X=\mu+h) \ge \Pr(X=\mu-h-1) \quad \text{dla } 0\le h\le \mu-1.
\]

\item[(ii)] Na tej podstawie wykaż, że 
\[
\Pr(X\ge \mu) \ge \tfrac{1}{2}.
\]
\end{enumerate}
\subsection*{Rozwiązanie}

Dla \( X\sim\mathrm{Poisson}(\mu) \) mamy
\[
\Pr(X=k)=e^{-\mu}\frac{\mu^k}{k!}.
\]

\subsubsection*{(i)}
Dla \(0\le h\le\mu-1\) rozważmy stosunek
\[
R:=\frac{\Pr(X=\mu+h)}{\Pr(X=\mu-h-1)}
=\frac{e^{-\mu}\mu^{\mu+h}/(\mu+h)!}{e^{-\mu}\mu^{\mu-h-1}/(\mu-h-1)!}
=\frac{\mu^{2h+1}}{(\mu-h)(\mu-h+1)\cdots(\mu+h)}.
\]

Można to zapisać jako
\[
R=\prod_{s=-h}^{h}\frac{\mu}{\mu+s}.
\]

Parując czynniki dla \(s\) i \(-s\) (dla \(s=1,\dots,h\)) oraz pozostawiając czynnik \(s=0\), otrzymujemy
\[
\frac{\mu}{\mu-s}\cdot\frac{\mu}{\mu+s}=\frac{\mu^2}{\mu^2-s^2}\ge1,
\]
bo \(\mu^2-s^2\le\mu^2\). Każda taka para jest więc nie mniejsza niż 1, a czynnik dla \(s=0\) to 1, zatem \(R\ge1\).

Stąd
\[
\Pr(X=\mu+h)\ge\Pr(X=\mu-h-1)
\]
dla każdego \(0\le h\le\mu-1\), co kończy część (i).

\subsubsection*{(ii)}
Z części (i) dla każdego \(h=0,1,\dots,\mu-1\) mamy 
\(\Pr(X=\mu+h)\ge\Pr(X=\mu-h-1)\). Sumując po \(h=0,\dots,\mu-1\) otrzymujemy
\[
\sum_{h=0}^{\mu-1}\Pr(X=\mu+h)\ge\sum_{h=0}^{\mu-1}\Pr(X=\mu-h-1).
\]

Lewa suma to część masy rozkładu od \(\mu\) wzwyż, prawa zaś – od 0 do \(\mu-1\). Zatem
\[
\Pr(X\ge\mu)\ge\Pr(X\le\mu-1).
\]
Ponieważ
\[
1=\Pr(X\ge\mu)+\Pr(X\le\mu-1),
\]
otrzymujemy
\[
\Pr(X\ge\mu)\ge\frac12.
\]
\subsection*{Odpowiedź}
\[
\boxed{\text{tutaj daj odpowiedz}}
\]
